\subsection{Integrable localization of the finite large-jump process}
For a zero-start c\`adl\`ag local martingale $N$, write
\[
 J_t=\sum_{0<s\leq t\wedge T,\,|\Delta N_s|>c}\Delta N_s,
 \qquad
 V_t=\sum_{0<s\leq t\wedge T,\,|\Delta N_s|>c}|\Delta N_s|.
\]
A c\`adl\`ag path has only finitely many jumps larger than $c>0$ on a
finite interval. Thus $J$ has finite variation and is constant after $T$.
Using the deterministic horizons
\[
 H_n=T+(n+1),
\]
define the strict variation passages
\[
 \rho_n=\min\{H_n,\inf\{s:V_s>n+1\}\}.
\]
The horizons are cofinal. Monotonicity of both the variation levels and
the horizons makes $(\rho_n)$ an increasing sequence of stopping times.
Finite total variation and constancy after $T$ imply, on every path, that
\[
 \rho_n=H_n
\]
for all sufficiently large $n$. Hence $(\rho_n)$ is an exhaustive localizing
sequence.

At first use the strict-prefix source
\[
 C^n_t=\begin{cases}
 J_t,&t<\rho_n,\\
 J_{\rho_n-},&t\geq\rho_n.
 \end{cases}
\]
It is adapted and c\`adl\`ag, with finite variation at most $n+1$ up to
$H_n$. Its terminal cumulative variation therefore belongs to $L^1$.
Each $C^n$ supplies normalized adapted c\`adl\`ag finite-variation $L^1$
data. The value-truncation theorem yields a predictable finite-variation
limit and a martingale residual. The source agrees exactly with $J$ on
the open prefix $t<\rho_n$.

Refine $\rho_n$ with a martingale localizer $\lambda_n$ and an
absolute-passage localizer $\kappa_n$, putting
$\tau_n=\rho_n\wedge\lambda_n\wedge\kappa_n$. These finite stops satisfy
\[
 \Var(J^{\tau_n};[0,H_n])\leq n+1+|\Delta J_{\tau_n}|.
\]
The next subsection proves integrability of the boundary term.
Thus every closed source $J^{\tau_n}$ also supplies finite-variation $L^1$
data and a value-truncation predictable projection. Fix all these choices
as one family. Each of $(\rho_n),(\lambda_n),(\kappa_n)$ is almost surely
increasing and exhaustive, so their minimum has the same properties.
For $m\leq n$, stopping absorption gives
\[
 (J^{\tau_n})^{\tau_m}\sim J^{\tau_m}.
\]
Here $\sim$ denotes indistinguishability: everywhere monotonicity of the
stops is not required.
For $m\leq n$, the difference between the upper predictable projection
stopped at $\tau_m$ and the lower projection is a local martingale, by this
closed-source identity. It is also strongly predictable, right-continuous,
and of finite variation. Rigidity of predictable finite-variation local
martingales makes it indistinguishable from its initial value. Both
projections start at zero almost surely, so their difference vanishes.
This establishes compatibility on ordered overlaps.
